% =====================================================================
%  ch05a.tex —— 第 5 章 阿贝尔簇 · §5.1
%  原书 p.573–580（自第 5 章标题 “Abelian varieties” 起，
%  止于 §5.2 标题 “Objects representing equivalences, and groups of
%  auto-equivalences” 之前；§5.2 起归 ch05b.tex）
% =====================================================================

% [原文 p.573]
\chapter{阿贝尔簇}
\section{阿贝尔簇上凝聚层范畴之间的等价}

\takeaway{Mukai 用 Poincaré 丛给出 $\Db(A)\cong\Db(\hat A)$。本节把由 $A\times\hat A$ 参数化的自等价 $\Phi_{(a,\alpha)}$ 汇成核 $S_A$，把任一等价化为同构 $f_\cE:A\times\hat A\to B\times\hat B$，并证明：两个凝聚层导出范畴等价，则 $A\times\hat A$ 与 $B\times\hat B$ 等距同构（定理 5.1.16）。}

本章研究阿贝尔簇上凝聚层的导出范畴，以及它们的自等价群。设 $A$ 是域 $k$ 上维数为 $n$ 的阿贝尔簇。记 $m:A\times A\to A$ 为合成态射（假定它定义在 $k$ 上），$e$ 为群结构的单位 $k$-点。对任意 $k$-点 $a\in A$，平移自同构 $m(\,\cdot\,,a):A\to A$ 记作 $T_a$。

记 $\hat A$ 为对偶阿贝尔簇，即 $A$ 上属于 $\Pic^0(A)$ 的线丛的模空间。此外，$\hat A$ 是细模空间，因此存在积 $A\times\hat A$ 上的泛线丛 $\cP$，称为 \emph{Poincaré 丛}。它由下述条件唯一确定：对任意 $k$-点 $\alpha\in\hat A$，$\cP$ 在 $A\times\{\alpha\}$ 上的限制同构于 $\Pic^0(A)$ 中对应于 $\alpha$ 的线丛；并且 $\cP|_{\{e\}\times\hat A}$ 平凡。

\yizhu{这里的“细”是关键：$\hat A$ 是细模空间，故存在泛线丛；再要求它在 $\{e\}\times\hat A$ 上平凡，才把 $\cP$ 唯一钉住——否则还可以乘上任一从 $\hat A$ 提升的线丛。命题 5.1.2 说相应的 $\Phi_\cP$ 是等价，这正是 Mukai 的原始结果，也是本节全部构造的起点。【来源：原书 p.573–574 与所引 [29]（Mukai，A 级）】}

\begin{definition}{5.1}
以下把 $A$ 上对应于 $k$-点 $\alpha\in\hat A$ 的线丛记作 $\cP_\alpha$。又，给定若干阿贝尔簇 $A_1,\dots,A_m$ 与 $k$-点 $(\alpha_1,\dots,\alpha_m)\in\hat A_1\times\cdots\times\hat A_m$，我们把积 $A_1\times\cdots\times A_k$ 上的线丛 $\cP_{\alpha_1}\boxtimes\cdots\boxtimes\cP_{\alpha_k}$ 记作 $\cP_{(\alpha_1,\dots,\alpha_k)}$。
\end{definition}

% [Tutor Note: 定义 5.1 前句用 $A_1,\dots,A_m$ 与 $(\alpha_1,\dots,\alpha_m)$，后句却写 $\cP_{(\alpha_1,\dots,\alpha_k)}$ 与 $A_1\times\cdots\times A_k$；两处下标应统一为 $m$。原书如此，照录。]

对阿贝尔簇的任意同态 $f:A\to B$，可以定义对偶同态 $\hat f:\hat B\to\hat A$。逐点来看，它把点 $\beta\in\hat B$ 送到 $\alpha\in\hat A$，当且仅当 $A$ 上的线丛 $f^*\cP_\beta$ 与 $\cP_\alpha$ 重合。

双对偶（二次对偶）阿贝尔簇 $\widehat{\hat A}$ 借助 $A\times\hat A$ 与 $\hat A\times\widehat{\hat A}$ 上的 Poincaré 丛自然地与 $A$ 等同。换言之，存在唯一同构 $\kappa_A:A\iso\widehat{\hat A}$，使得 Poincaré 丛 $P_{\hat A}$ 在同构 $1\times\kappa_A:\hat A\times A\iso\hat A\times\widehat{\hat A}$ 下的拉回与 Poincaré 丛 $\cP_A$ 重合，即 $(1\times\kappa_A)^*P_{\hat A}\cong\cP_A$。于是，$\widehat{\phantom{x}}$ 是阿贝尔簇范畴上的对合：它是反变函子，其平方同构于恒等函子，即 $\kappa:\id\iso\widehat{\widehat{\phantom{x}}}$。

% [Tutor Note: 本段页图把 $\hat A\times\widehat{\hat A}$ 上的 Poincaré 丛印作斜体 $P_{\hat A}$，而同一句的 $\cP_A$ 用花体 $\cP$；同一丛出现两种字体，照录。]

Poincaré 丛 $\cP$ 给出 $A$ 与 $\hat A$ 上凝聚层导出范畴之间正合等价的一个例子，而这两个簇一般并不彼此同构。考虑投影
\begin{equation*}
A\xleftarrow{\ p\ }A\times\hat A\xrightarrow{\ q\ }\hat A
\end{equation*}
以及由式 (7) 定义的函子 $\Phi_\cP:\Db(A)\to\Db(\hat A)$：
\begin{equation*}
\Phi_\cP(\,\cdot\,)=\Rf q_*\bigl(\cP\otimes p^*(\,\cdot\,)\bigr).
\end{equation*}

\yizhu{$A$ 与 $\hat A$ 一般并不同构，它们的导出范畴却等价：这是历史上第一个“不同簇、同范畴”的例子（见引言）。$\Phi_\cP$ 就是 Fourier--Mukai 变换的原型，$A\times\hat A$ 上的 Poincaré 丛是它的核；命题 5.1.2 中的 $\Psi_\cP$ 则是它的逆（差一个平移与群逆）。【来源：原书引言（p.513）与 §5.1（p.573–574），并见所引 [29]（Mukai，A 级）】}

下述命题由 [29] 证明。

% [原文 p.574]
\begin{proposition}[29]{5.1.2}
设 $\cP$ 是 $A\times\hat A$ 上的 Poincaré 丛。则函子 $\Phi_\cP:\Db(A)\to\Db(\hat A)$ 是正合等价，并且存在函子同构
\begin{equation*}
\Psi_\cP\circ\Phi_\cP\cong(-1_A)^*[n],
\end{equation*}
其中 $(-1_A)$ 是 $A$ 的群逆。
\end{proposition}

\begin{remark}{5.1.3}
在 [29] 中，这一断言是对代数闭域上的阿贝尔簇证明的。然而它对任意域也成立，因为对偶簇与 Poincaré 丛总定义在同一个域上（例如见 [32]）。至于范畴等价这一断言，则由引理 5.1.9 推出。
\end{remark}

考虑 $k$-点 $(a,\alpha)\in A\times\hat A$。对每个这样的点，都可以按如下规则给它配一个 $\Db(A)$ 到自身的函子：
\begin{equation}
\Phi_{(a,\alpha)}(\,\cdot\,):=T_{a*}(\,\cdot\,)\otimes\cP_\alpha=T^*_{-a}(\,\cdot\,)\otimes\cP_\alpha. \tag{49}
\end{equation}

函子 $\Phi_{(a,\alpha)}$ 由积 $A\times A$ 上的层
\begin{equation}
S_{(a,\alpha)}=\cO_{\Gamma_a}\otimes p^*_2(\cP_\alpha) \tag{50}
\end{equation}
代表，其中 $\Gamma_a$ 表示平移自同构 $T_a$ 的图。函子 $\Phi_{(a,\alpha)}$ 显然是自等价。

由 $A\times\hat A$ 参数化的全部函子 $\Phi_{(a,\alpha)}$ 可以汇成单个函子，从 $\Db(A\times\hat A)$ 到 $\Db(A\times A)$，它把点的结构层 $\cO_{(a,\alpha)}$ 送到 $S_{(a,\alpha)}$。（注意，这一条件并不唯一确定该函子，而只把它确定到“乘以从 $A\times\hat A$ 提升的线丛”为止。）

我们把所需的函子 $\Phi_{S_A}:\Db(A\times\hat A)\to\Db(A\times A)$ 定义为另外两个函子的复合。考虑对象
\begin{equation*}
P_A=p^*_{14}\cO_\Delta\otimes p^*_{23}\cP\in\Db\bigl((A\times\hat A)\times(A\times A)\bigr),
\end{equation*}
并记 $\mu_A:A\times A\to A\times A$ 为把 $(a_1,a_2)$ 送到 $(a_1,m(a_1,a_2))$ 的态射。我们得到两个函子
\begin{equation*}
\Phi_{P_A}:\Db(A\times\hat A)\to\Db(A\times A),\qquad
\Rf\mu_{A*}:\Db(A\times A)\to\Db(A\times A).
\end{equation*}

\begin{definition}{5.1.4}
函子 $\Phi_{S_A}$ 是复合 $\Rf\mu_{A*}\circ\Phi_{P_A}$。
\end{definition}

\begin{proposition}{5.1.5}
函子 $\Phi_{S_A}$ 是范畴等价。对任意 $k$-点 $(a,\alpha)\in A\times\hat A$，它把
\begin{enumerate}[label=\alph*),leftmargin=2.8em,itemsep=1pt,topsep=2pt]
\item 点 $(a,\alpha)$ 的结构层 $\cO_{(a,\alpha)}$ 送到由 (50) 定义的层 $S_{(a,\alpha)}$，
\item $A\times\hat A$ 上的线丛 $\cP_{(\alpha,a)}$ 送到对象 $\cO_{\{-a\}\times A}\otimes p^*_2\cP_\alpha[n]$。
\end{enumerate}
\end{proposition}

\begin{proof}
按定义，$\Phi_{S_A}$ 是函子 $\Rf\mu_{A*}$ 与 $\Phi_{P_A}$ 的复合，二者都是等价：前者显然；后者由命题 2.1.7 与命题 5.1.2 推出。

函子 $\Phi_{P_A}$ 把点的结构层 $\cO_{(a,\alpha)}$ 送到 $\cO_{A\times\{a\}}\otimes p^*_1\cP_\alpha$。而 $\Rf\mu_{A*}$ 把 $\cO_{A\times\{a\}}\otimes p^*_1\cP_\alpha$ 送到 $\cO_{\Gamma_a}\otimes p^*_1(\cP_\alpha)$。

% [Tutor Note: 页图此处印 $p^*_1(\cP_\alpha)$，而 (50) 与命题 5.1.5 b) 均印 $p^*_2$；按 (50) 应为 $p^*_2$。原书如此，照录。]
% [原文 p.575]
同法，应用命题 5.1.2 可知，$\Phi_{P_A}$ 把线丛 $\cP_{(\alpha,a)}$ 送到对象 $\cO_{\{-a\}\times A}\otimes p^*_2\cP_\alpha[n]$，而 $\Rf\mu_{A*}$ 把 $\cO_{\{-a\}\times A}\otimes p^*_2\cP_\alpha[n]$ 送到它自身。
\end{proof}

设 $A$ 与 $B$ 是两个凝聚层导出范畴等价的阿贝尔簇。固定其中一个等价。由定理 3.2.2，它可由积上的一个对象代表。于是存在对象 $\cE\in\Db(A\times B)$ 与等价 $\Phi_\cE:\Db(A)\iso\Db(B)$。考虑由 (11) 定义的函子
\begin{equation*}
\mathrm{Ad}_\cE:\Db(A\times A)\iso\Db(B\times B),
\end{equation*}
它是等价；再考虑函子的复合 $\Phi^{-1}_{S_B}\circ\mathrm{Ad}_\cE\circ\Phi_{S_A}$。

\begin{definition}{5.1.6}
我们把代表函子 $\Phi^{-1}_{S_B}\circ\mathrm{Ad}_\cE\circ\Phi_{S_A}$ 的对象记作 $\cJ(\cE)$。
\end{definition}

于是有交换图
\begin{equation}
\begin{tikzcd}[row sep=2.6em,column sep=3.2em,ampersand replacement=\&]
\Db(A\times\hat A) \arrow[r, "\Phi_{S_A}"] \arrow[d, "\Phi_{\cJ(\cE)}"]
  \& \Db(A\times A) \arrow[d, "\mathrm{Ad}_{\cE}"] \\
\Db(B\times\hat B) \arrow[r, "\Phi_{S_B}"]
  \& \Db(B\times B)
\end{tikzcd}
\tag{51}
\end{equation}

下述定理使我们可以算出对象 $\cJ(\cE)$；它是描述具有等价凝聚层导出范畴的阿贝尔簇的主要工具。

\begin{theorem}{5.1.7}
存在阿贝尔簇的同态 $f_\cE:A\times\hat A\to B\times\hat B$（它是同构），以及 $A\times\hat A$ 上的线丛 $\cL_\cE$，使得对象 $\cJ(\cE)$ 同构于 $i_*(\cL_\cE)$，其中 $i$ 是 $A\times\hat A$ 作为同构 $f_\cE$ 的图到 $(A\times\hat A)\times(B\times\hat B)$ 的嵌入。
\end{theorem}

在证明这一定理之前，我们先陈述两个引理，它们使我们能够假定域 $k$ 代数闭。记 $\bar k$ 为 $k$ 的代数闭包，置 $\bar X:=X\times_{\Spec(k)}\Spec(\bar k)$，并记 $\bar\cF$ 为 $\cF$ 在态射 $\bar X\to X$ 下的逆像。

\begin{lemma}[37]{5.1.8}
设 $\cF$ 是光滑簇 $X$ 上的凝聚层。假设存在闭子簇 $j:Z\into\bar X$ 与 $Z$ 上的可逆层 $\cL$，使得 $\bar\cF\cong j_*\cL$。则存在闭子簇 $i:Y\into X$ 与 $Y$ 上的可逆层 $\cM$，使得 $\cF\cong i_*\cM$ 且 $j=\overline{\subset}$。
\end{lemma}

% [Tutor Note: 引理 5.1.8 末字页图印作带上划线的 $\subset$（即 $j=\overline{\subset}$），当是 $j=\bar\iota$ 之误排，其中 $\iota$ 为 $i:Y\into X$ 在 $\bar k$ 上的基变换；语义为 $j$ 是 $i$ 的基变换。原书如此，照录。]

下一个引理告诉我们：函子全忠实（或为等价）这一性质在域扩张下稳定。

\yizhu{两个引理合起来是一次下降论证：引理 5.1.8 把“闭子簇上的可逆层”沿域扩张 $\bar k/k$ 下降回 $k$，引理 5.1.9 说明全忠实与等价在域扩张下不变。于是定理 5.1.7 的证明可设 $k$ 代数闭，第 1–3 步的几何论证才得以展开；否则摩天楼层 $\cO_e$、仿射邻域与极大理想 $\mathfrak m$ 这些工具都要换成一般域上的说法。【来源：原书 p.575–576 与所引 [37]（A 级）】}

% [原文 p.576]
\begin{lemma}[37]{5.1.9}
设 $X$ 与 $Y$ 是域 $k$ 上的光滑射影簇，$\cE$ 是导出范畴 $\Db(X\times Y)$ 的对象。考虑域扩张 $F/k$ 与簇
\begin{equation*}
X'=X\times_{\Spec(k)}\Spec(F),\qquad Y'=Y\times_{\Spec(k)}\Spec(F).
\end{equation*}
设 $\cE'$ 是 $\cE$ 到范畴 $\Db(X'\times Y')$ 的提升。则函子 $\Phi_\cE:\Db(X)\to\Db(Y)$ 全忠实（或为等价），当且仅当 $\Phi_{\cE'}:\Db(X')\to\Db(Y')$ 全忠实（相应地，为等价）。
\end{lemma}

\begin{proof}[定理 5.1.7 的证明]
利用引理 5.1.8 与引理 5.1.9，我们可以过渡到域 $k$ 的代数闭包。

\emph{第 1 步。}记 $e\in A\times\hat A$ 与 $e'\in B\times\hat B$ 为群结构单位元所对应的闭点。考虑摩天层 $\cO_e$，并计算它在函子 $\Phi_{\cJ(\cE)}$ 下的像。按定义，
\begin{equation*}
\Phi_{\cJ(\cE)}=\Phi^{-1}_{S_B}\circ\mathrm{Ad}_\cE\circ\Phi_{S_A}.
\end{equation*}
由命题 5.1.5，函子 $\Phi_{S_A}$ 把 $\cO_e$ 送到 $A\times A$ 中对角线的结构层 $\cO_{\Delta(A)}$。对角线的结构层代表恒等函子，故由 (12) 可知 $\mathrm{Ad}_\cE(\cO_{\Delta(A)})$ 是 $B\times B$ 中对角线的结构层 $\cO_{\Delta(B)}$。再由命题 5.1.5，这个层在函子 $\Phi^{-1}_{S_B}$ 的作用下送到结构层 $\cO_{e'}$。

\emph{第 2 步。}于是我们看到
\begin{equation*}
\cJ(\cE)\otL\cO_{\{e\}\times(B\times\hat B)}\cong\cO_{\{e\}\times\{e'\}}.
\end{equation*}
由此可知，存在 $e$ 在 Zariski 拓扑下的仿射邻域 $U=\Spec(R)$，使得对象 $\cJ':=\cJ(\cE)|_{U\times(B\times\hat B)}$ 是凝聚层，其支撑与纤维 $\{e\}\times(B\times\hat B)$ 交于点 $\{e\}\times\{e'\}$。我们回忆：任意凝聚层的支撑都是闭子集。

现取点 $e'$ 在 $B\times\hat B$ 中的某个仿射邻域 $V=\Spec(S)$。$\cJ'$ 的支撑与补集 $B\times\hat B\setminus V$ 的交是闭子集，它到 $A\times\hat A$ 的投影是不含点 $e$ 的闭子集。于是，必要时缩小 $U$，可以假定它仍然是仿射的，并且 $\cJ'$ 的支撑含于 $U\times V$。这意味着存在 $U\times V$ 上的凝聚层 $\cF$，使得 $j_*(\cF)=\cJ'$，其中 $j$ 是 $U\times V$ 到 $U\times(B\times\hat B)$ 的嵌入。记 $\cM$ 为与层 $\cF$ 对应的有限生成 $R\otimes S$-模，即 $\cF=\tilde{\cM}$。此外我们注意，$\cM$ 作为 $R$-模也是有限生成的，因为凝聚层 $\cJ'=j_*\cF$ 沿投影的直接像是凝聚层。

设 $\mathfrak m$ 是 $R$ 中对应于点 $e$ 的极大理想。我们知道，
\begin{equation*}
\cM\otimes_R R/\mathfrak m\cong R/\mathfrak m.
\end{equation*}
于是存在 $R$-模同态 $\varphi:R\to\cM$，它在与 $R/\mathfrak m$ 张量之后成为同构。因此，凝聚层 $\Ker\varphi$ 与 $\Coker\varphi$ 的支撑都不含点 $e$。所以，
% [原文 p.577]
把 $U$ 换成与层 $\Ker\varphi$、$\Coker\varphi$ 的支撑都不交的、$e$ 的某个更小仿射邻域，便知 $\varphi$ 是同构。因此存在子概形 $X(U)\subset U\times(B\times\hat B)$，使得投影 $X(U)\to U$ 是同构，且
\begin{equation*}
\cJ'=\cJ(\cE)|_{U\times(B\times\hat B)}\cong\cO_{X(U)}.
\end{equation*}

\emph{第 3 步。}至此我们已证明：对任意闭点 $(a,\alpha)\in U$，
\begin{equation*}
\Phi_{\cJ(\cE)}(\cO_{(a,\alpha)})\cong\cO_{(b,\beta)}
\end{equation*}
其中 $(b,\beta)$ 是 $B\times\hat B$ 的某个闭点。现在任取闭点 $(a,\alpha)\in A\times\hat A$，我们总可以把它写成和 $(a,a')=(a_1,\alpha_1)+(a_2,\alpha_2)$，其中点 $(a_1,\alpha_1)$ 与 $(a_2,\alpha_2)$ 属于 $U$。记 $(b_1,\beta_1)$ 与 $(b_2,\beta_2)$ 为这些点在函子 $\Phi_{\cJ(\cE)}$ 下的像。我们知道，函子 $\Phi_{S_A}$ 把结构层 $\cO_{(a,\alpha)}$ 送到层 $S_{(a,\alpha)}$。记 $\cG$ 为对象 $\mathrm{Ad}_\cE(S_{(a,\alpha)})$。我们用 (12) 来计算它。有同构
\begin{equation*}
\Phi_\cG\cong\Phi_\cE\circ\Phi_{(a,\alpha)}\circ\Phi_\cE^{-1}.
\end{equation*}

然而，函子 $\Phi_{(a,\alpha)}$ 按定义 (49) 等于 $T^*_a(\,\cdot\,)\otimes\cP_\alpha$，并可表为复合 $\Phi_{(a_1,\alpha_1)}\Phi_{(a_2,\alpha_2)}$。于是我们得到一串同构
\begin{equation*}
\begin{aligned}
\Phi_\cG&\cong\Phi_\cE\circ\Phi_{(a,\alpha)}\circ\Phi_\cE^{-1}
\cong\Phi_\cE\circ\Phi_{(a_1,\alpha_1)}\circ\Phi_\cE^{-1}\\
&\cong\Phi_\cE\circ\Phi_{(a_2,\alpha_2)}\circ\Phi_\cE^{-1}
\cong\Phi_{(b_1,\beta_1)}\circ\Phi_{(b_2,\beta_2)}\cong\Phi_{(b,\beta)},
\end{aligned}
\end{equation*}
其中 $(b,\beta)=(b_1,\beta_1)+(b_2,\beta_2)$。因此，对象 $\cG$ 同构于 $S_{(b,\beta)}$。最终我们得到
\begin{equation*}
\Phi_{\cJ(\cE)}(\cO_{(a,\alpha)})\cong\cO_{(b,\beta)}
\end{equation*}
对任意闭点 $(a,\alpha)\in A\times\hat A$ 成立。

% [Tutor Note: 第 3 步两处把 $(a,\alpha)$ 印作 $(a,a')$：一处在“写成和 $(a,a')=(a_1,\alpha_1)+(a_2,\alpha_2)$”，一处在下段“对任意闭点 $(a,a')$”。按上下文应为 $(a,\alpha)$。原书如此，照录。]
% [Tutor Note: 上条同构链中第二、三个 $\cong$ 各只保留复合的一个因子；按 $\Phi_{(a,\alpha)}=\Phi_{(a_1,\alpha_1)}\Phi_{(a_2,\alpha_2)}$，前两步应合写为 $\Phi_\cE\circ\Phi_{(a_1,\alpha_1)}\Phi_{(a_2,\alpha_2)}\circ\Phi_\cE^{-1}$。原书如此，照录。]
% [Tutor Note: 原书说 $\Phi_{(a,\alpha)}$ “按定义 (49) 等于 $T^*_a(\,\cdot\,)\otimes\cP_\alpha$”，但 (49) 印的是 $T^*_{-a}$；且 $T^*_a=T_{-a*}$，与 (49) 的 $T_{a*}$ 并不同。此处应为 $T^*_{-a}$。原书如此，照录。]

现在重复第 2 步的做法：对任意闭点 $(a,a')$，都能找到邻域 $W$ 与子概形 $X(W)\subset W\times(B\times\hat B)$，使得投影 $X(W)\to W$ 是同构，且 $\cJ|_{W\times(B\times\hat B)}\cong\cO_{X(W)}$。把这些邻域粘起来，便得到子簇 $i:X\into(A\times\hat A)\times(B\times\hat B)$，使得投影 $X\to A\times\hat A$ 是同构，且层 $\cJ(\cE)$ 同构于 $i_*\cL$，其中 $\cL$ 是 $X$ 上的线丛。子簇 $X$ 定义从 $A\times\hat A$ 到 $B\times\hat B$ 的同态，它诱导导出范畴的等价。因此，这个同态是同构。

特别地，由此立即推出：若两个阿贝尔簇 $A$ 与 $B$ 的凝聚层导出范畴等价，则簇 $A\times\hat A$ 与 $B\times\hat B$ 同构。下面我们将说明，这个同构还必须满足某个附加条件（见命题 5.1.15）。
\end{proof}

% [原文 p.578]
\begin{corollary}{5.1.10}
同构 $f_\cE$ 把 $k$-点 $(a,\alpha)\in A\times\hat A$ 送到点 $(b,\beta)\in B\times\hat B$，当且仅当由公式 (49) 定义的等价
\begin{equation*}
\Phi_{(a,\alpha)}:\Db(A)\iso\Db(A),\qquad
\Phi_{(b,\beta)}:\Db(B)\iso\Db(B)
\end{equation*}
满足如下关系：
\begin{equation*}
\Phi_{(b,\beta)}\circ\Phi_\cE\cong\Phi_\cE\circ\Phi_{(a,\alpha)},
\end{equation*}
或者，用对象来写，
\begin{equation*}
T_{b*}\cE\otimes\cP_\beta\cong T_{-a*}\cE\otimes\cP_\alpha=T^*_a\cE\otimes\cP_\alpha.
\end{equation*}
\end{corollary}

\begin{proof}
由定理 5.1.7，$\Phi_{\cJ(\cE)}$ 把 $(a,\alpha)$ 的结构层 $\cO_{(a,\alpha)}$ 送到 $(b,\beta)=f_\cE(a,\alpha)$ 的结构层 $\cO_{(b,\beta)}$。由命题 5.1.5，$\Phi_{S_A}$ 把 $\cO_{(a,\alpha)}$ 送到 $S_{(a,\alpha)}$。而层 $S_{(a,\alpha)}$ 代表函子 $\Phi_{(a,\alpha)}=T_{a*}(\,\cdot\,)\otimes\cP_\alpha$。

现在利用图 (51) 可以看出，$f_\cE$ 把 $(a,\alpha)$ 送到 $(b,\beta)$，当且仅当 $S_{(b,\beta)}\cong\mathrm{Ad}_\cE(S_{(a,\alpha)})$。应用公式 (12) 便得 $\Phi_{(b,\beta)}\cong\Phi_\cE\circ\Phi_{(a,\alpha)}\circ\Phi_\cE^{-1}$。
\end{proof}

下文需要 $\cJ(\cE)$ 在如下特殊情形下的显式公式：$A=B$，且等价 $\Phi_\cE$ 等于由公式 (49) 定义的 $\Phi_{(a,\alpha)}$。

\begin{proposition}{5.1.11}
设 $A=B$。考虑 $A\times A$ 上代表由 (49) 给出的等价 $\Phi_{(a,\alpha)}$ 的对象 $S_{(a,\alpha)}$。则 $\cJ(S_{(a,\alpha)})$ 等于 $\Delta_*\cP_{(\alpha,-a)}$，其中 $\Delta$ 是 $A\times\hat A$ 到 $(A\times\hat A)\times(A\times\hat A)$ 的对角嵌入，而 $\cP_{(\alpha,a)}$ 是 5.1.1 中定义的 $A\times\hat A$ 上的线丛。
\end{proposition}

% [Tutor Note: 命题 5.1.11 陈述中前一处印 $\cP_{(\alpha,-a)}$，后一处却印 $\cP_{(\alpha,a)}$；又“5.1.1”应是“定义 5.1”。两处下标与编号均疑为原书笔误，照录。]

\begin{proof}
由命题 5.1.5，$\Phi_{S_A}$ 把 $\cO_{(a',\alpha')}$ 送到 $A\times A$ 上的层 $S_{(a',\alpha')}$（见 (50)）。此外，$\mathrm{Ad}_{S_{(a,\alpha)}}$ 把 $S_{(a',\alpha')}$ 送到自身，因为由公式 (12)，对象 $\mathrm{Ad}_{S_{(a,\alpha)}}(S_{(a',\alpha')})$ 代表函子
\begin{equation*}
\Phi_{(a,\alpha)}\circ\Phi_{(a',\alpha')}\circ\Phi^{-1}_{(a,\alpha)},
\end{equation*}
而它同构于 $\Phi_{(a',\alpha')}$，因为这类函子两两交换。于是我们看到，由 $\cJ(S_{(a,\alpha)})$ 定义的函子把每个点的结构层送到自身；从而层 $\cJ(S_{(a,\alpha)})$ 是某个集中在对角线上的线丛 $L$。

现在要找出线丛 $L$，可以追问该函子把丛 $\cP_{(\alpha',a')}$ 送到何处。再次应用命题 5.1.5 可知，函子 $\Phi_{S_A}$ 把 $\cP_{(\alpha',a')}$ 送到对象 $\cO_{\{-a'\}\times A}\otimes p^*_2(\cP_{\alpha'})[n]$。其次容易看出，在函子 $\mathrm{Ad}_{S_{(a,\alpha)}}$ 的作用下它送到对象 $\cO_{\{-a'+a\}\times A}\otimes p^*_2(\cP_{\alpha'+\alpha})[n]$。因此，在由层 $\cJ(S_{(a,\alpha)})$ 给出的函子作用下，丛 $\cP_{(\alpha',a')}$ 送到丛 $\cP_{(\alpha'+\alpha,a'-a)}$。也就是说，$L$ 同构于 $\cP_{(\alpha,-a)}$。
\end{proof}

% [原文 p.579]
对阿贝尔簇 $A$ 与 $B$，记 $\mathrm{Eq}(A,B)$ 为 $\Db(A)$ 到 $\Db(B)$ 的全部正合等价（模同构）之集。我们引入两个群胚 $\mathfrak A$ 与 $\mathfrak D$（即所有态射都可逆的范畴）。二者的对象都是阿贝尔簇。$\mathfrak A$ 的态射是阿贝尔簇作为代数群之间的同构。$\mathfrak D$ 的态射是阿贝尔簇上凝聚层导出范畴之间的正合等价；也就是说，
\begin{equation*}
\Mor_{\mathfrak A}(A,B):=\mathrm{Iso}(A,B)
\qquad\text{与}\qquad
\Mor_{\mathfrak D}(A,B):=\mathrm{Eq}(A,B).
\end{equation*}

定理 5.1.7 给出从集合 $\mathrm{Eq}(A,B)$ 到集合 $\mathrm{Iso}(A\times\hat A,B\times\hat B)$ 的映射，它把等价 $\Phi_\cE$ 送到同构 $f_\cE$。我们考虑从 $\mathfrak D$ 到 $\mathfrak A$ 的映射 $F$：它把阿贝尔簇 $A$ 送到簇 $A\times\hat A$，并按上述方式作用在态射上。

\begin{proposition}{5.1.12}
映射 $F:\mathfrak D\to\mathfrak A$ 是函子。
\end{proposition}

\begin{proof}
为证明这一断言，只需说明 $F$ 保持态射的复合。考虑三个阿贝尔簇 $A$、$B$、$C$。设 $\cE$ 与 $\cF$ 分别是 $\Db(A\times B)$ 与 $\Db(B\times C)$ 的对象，使得函子
\begin{equation*}
\Phi_\cE:\Db(A)\to\Db(B)\qquad\text{与}\qquad \Phi_\cF:\Db(B)\to\Db(C)
\end{equation*}
都是等价。记 $\cG$ 为 $\Db(A\times C)$ 中代表这两个函子之复合的对象。

关系 (10) 给出同构 $\mathrm{Ad}_\cG\cong\mathrm{Ad}_\cF\circ\mathrm{Ad}_\cE$。于是我们看到
\begin{equation*}
\Phi_{\cJ(\cF)}\circ\Phi_{\cJ(\cE)}
\cong\bigl(\Phi^{-1}_{S_A}\circ\mathrm{Ad}_\cF\circ\Phi_{S_A}\bigr)
\circ\bigl(\Phi^{-1}_{S_A}\circ\mathrm{Ad}_\cE\circ\Phi_{S_A}\bigr)
\cong\Phi^{-1}_{S_A}\circ\mathrm{Ad}_\cG\circ\Phi_{S_A}
\cong\Phi_{\cJ(\cG)}.
\end{equation*}
由定理 5.1.7，对象 $\cJ(\cE)$、$\cJ(\cF)$、$\cJ(\cG)$ 都是分别集中在同构 $f_\cE$、$f_\cF$、$f_\cG$ 的图上的线丛。于是我们得到关系 $f_\cG=f_\cF\cdot f_\cE$。
\end{proof}

\begin{corollary}{5.1.13}
设 $A$ 是阿贝尔簇，$\Phi_\cE$ 是导出范畴 $\Db(A)$ 的自等价。则对应 $\Phi_\cE\mapsto f_\cE$ 定义群同态
\begin{equation*}
\gamma_A:\Auteq\Db(A)\to\Aut(A\times\hat A).
\end{equation*}
\end{corollary}

于是存在函子 $F:\mathfrak D\to\mathfrak A$。我们的下一个目标是描述这个函子。为此必须确定：$\mathrm{Iso}(A\times\hat A,B\times\hat B)$ 中哪些元素可以对某个 $\cE$ 实现为 $f_\cE$；还要回答两个等价 $\cE_1$、$\cE_2$ 何时满足 $f_{\cE_1}=f_{\cE_2}$。

考虑任意态射 $f:A\times\hat A\to B\times\hat B$。把它写成矩阵
\begin{equation*}
f=\begin{pmatrix}\alpha&\beta\\ \gamma&\delta\end{pmatrix}
\end{equation*}
比较方便，其中 $\alpha$ 是从 $A$ 到 $B$ 的态射，$\beta$ 从 $\hat A$ 到 $B$，$\gamma$ 从 $A$ 到 $\hat B$，$\delta$ 从 $\hat A$ 到 $\hat B$。每个态射 $f$ 定义另外两个从 $B\times\hat B$ 到 $A\times\hat A$ 的态射 $\hat f$ 与 $\tilde f$，它们有下述矩阵形式：
\begin{equation*}
\hat f=\begin{pmatrix}\hat\delta&\hat\beta\\ \hat\gamma&\hat\alpha\end{pmatrix}
\qquad\text{与}\qquad
\tilde f=\begin{pmatrix}\hat\delta&-\hat\beta\\ -\hat\gamma&\hat\alpha\end{pmatrix}.
\end{equation*}

% [原文 p.580]
我们定义集合 $U(A\times\hat A,B\times\hat B)$ 为 $\mathrm{Iso}(A\times\hat A,B\times\hat B)$ 中满足“$\tilde f$ 与 $f$ 的逆重合”的那些 $f$ 构成的子集，即
\begin{equation*}
U(A\times\hat A,B\times\hat B):=\bigl\{\,f\in\mathrm{Iso}(A\times\hat A,B\times\hat B)\bigm|\tilde f=f^{-1}\,\bigr\}.
\end{equation*}
若 $B=A$，则把这个集合记作 $U(A\times\hat A)$。我们注意，$U(A\times\hat A)$ 是 $\Aut(A\times\hat A)$ 的子群。

\begin{definition}{5.1.14}
若同构 $f:A\times\hat A\iso B\times\hat B$ 属于 $U(A\times\hat A,B\times\hat B)$，则称它是\emph{等距}的（isometric）。
\end{definition}

\yizhu{“等距”一词来自 $A\times\hat A$ 上由 Poincaré 丛 $\cP$ 给出的结构：条件 $\tilde f=f^{-1}$ 把 $f$ 的矩阵与它的对偶矩阵（转置并变号）绑在一起，正是“保持这一结构”的写法。命题 5.1.15 表明，任何导出等价给出的 $f_\cE$ 都自动等距，因此定理 5.1.16 所得的同构比一般的簇同构强。【来源：原书 p.579–580（定义 5.1.14、命题 5.1.15）与所引 [38]（A 级）】}

\begin{proposition}{5.1.15}
对任意等价 $\Phi_\cE:\Db(A)\iso\Db(B)$，同构 $f_\cE$ 是等距的。
\end{proposition}

\begin{proof}
必要时过渡到代数闭包，我们可以假定 $k$ 代数闭。为验证等式 $\tilde f_\cE=f_\cE^{-1}$，只需说明这两个态射在闭点上一致。设 $f_\cE$ 把 $(a,\alpha)\in A\times\hat A$ 送到 $(b,\beta)\in B\times\hat B$。我们必须证明 $\tilde f_\cE(b,\beta)=(a,\alpha)$，或者等价地，$\hat f_\cE(-b,\beta)=(-a,\alpha)$。

同构 $f_\cE$ 由阿贝尔子簇 $X\into A\times\hat A\times B\times\hat B$ 给出。因此我们必须证明
\begin{equation*}
\cP_{(0,0,\beta,-b)}\otimes\cO_X\cong\cP_{(\alpha,-a,0,0)}\otimes\cO_X,
\end{equation*}
或者等价地，
\begin{equation*}
\cJ':=\cP_{(-\alpha,a,\beta,-b)}\otimes\cJ(\cE)
\end{equation*}
同构于层 $\cJ(\cE)$。

由命题 5.1.11，$\cJ'$ 所给出的函子是对象 $\cJ(S_{(-a,-\alpha)})$、$\cJ(\cE)$、$\cJ(S_{(b,\beta)})$ 所代表函子的复合。因此 $\cJ'$ 与 $\cJ(\cE')$ 重合，其中 $\cE'$ 是 $\Db(A\times B)$ 中代表函子
\begin{equation*}
\Phi_{(b,\beta)}\circ\Phi_\cE\circ\Phi_{(-a,-\alpha)}
\end{equation*}
的对象。由推论 5.1.10，这个复合同构于函子 $\Phi_\cE$。这意味着对象 $\cE'$ 同构于 $\cE$，从而 $\cJ'=\cJ(\cE')\cong\cJ(\cE)$。
\end{proof}

作为定理 5.1.7 与命题 5.1.15 的推论，我们得到下述结果。

\begin{theorem}{5.1.16}
设 $A$ 与 $B$ 是域 $k$ 上的两个阿贝尔簇。若凝聚层导出范畴 $\Db(A)$ 与 $\Db(B)$ 作为三角范畴等价，则 $A\times\hat A$ 与 $B\times\hat B$ 之间存在等距同构。
\end{theorem}

其逆命题对代数闭域上特征 $0$ 的阿贝尔簇成立，[38] 已证。我们在 §5.3 中将给出这一事实的另一个证明。

\begin{corollary}{5.1.17}
对任意阿贝尔簇 $A$，只有有限多个互不同构的阿贝尔簇，其凝聚层导出范畴与 $\Db(A)$（作为三角范畴）等价。
\end{corollary}

\begin{proof}
[26] 证明了：对任意阿贝尔簇 $Z$，只有有限多个阿贝尔簇（模同构）能作为阿贝尔子簇嵌入 $Z$。把这一断言用于 $Z=A\times\hat A$，并利用定理 5.1.16，即得所需结果。
\end{proof}
